\documentclass[11pt,a4paper]{altacv}

\geometry{left=1.5cm,right=1.5cm,top=1.cm,bottom=1.2cm,columnsep=1.5cm}

\usepackage[utf8]{inputenc}
\usepackage[T1]{fontenc}
\usepackage[french]{babel}
\usepackage{paracol}
\usepackage{xcolor}
\usepackage{hyperref}
\usepackage{fontawesome5}
\usepackage{setspace}

\definecolor{SlateGrey}{HTML}{2E2E2E}
\definecolor{LightGrey}{HTML}{666666}
\definecolor{DarkPastelRed}{HTML}{8AC0D9}
\definecolor{PastelRed}{HTML}{00B0DA}
\definecolor{GoldenEarth}{HTML}{B4AD0C}
\colorlet{name}{black}
\colorlet{tagline}{PastelRed}
\colorlet{heading}{DarkPastelRed}
\colorlet{headingrule}{GoldenEarth}
\colorlet{subheading}{PastelRed}
\colorlet{accent}{PastelRed}
\colorlet{emphasis}{SlateGrey}
\colorlet{body}{LightGrey}

\renewcommand{\familydefault}{\sfdefault}

\name{\Large Guillaume BOILEAU}
\tagline{Chef de projet technique DSI | SI, build/run, risques, cycle en V, Agile \& interfaces techniques}
\personalinfo{
  \email{guillaume.boileaupro@gmail.com}
  \phone{+33 6 59 94 85 84}
  \location{Alpes-Maritimes, France}
  \linkedin{guillaume-boileau-phd-891b81265}
}

\setlength{\parskip}{0.75\baselineskip}

\begin{document}
\makecvheader
\vspace{-0.6cm}

\cvsection{Lettre de motivation}
\vspace{-0.4cm}

{\color{accent}\textbf{Objet :}} \textit{Candidature -- Chef de projet stratégique informatique -- Valbonne}

{\color{body}

Madame, Monsieur,

Je vous adresse ma candidature pour le poste de \textbf{Chef de projet stratégique informatique}. Votre offre m'intéresse particulièrement car elle combine des dimensions qui correspondent fortement à mon parcours : coordination de sujets techniques complexes, qualification de besoins, interface entre équipes, cycle de vie projet, gestion des risques, reporting, suivi d'actions et compréhension des enjeux build/run.

Mon parcours s'est construit à l'interface entre \textbf{développement logiciel, systèmes complexes, validation, analyse technique, documentation et coordination transverse}. J'ai travaillé dans des environnements où il est essentiel de comprendre rapidement les besoins, d'identifier les contraintes, de structurer les dépendances, de fédérer plusieurs parties prenantes et de produire des éléments de décision clairs.

Au \textbf{CNES}, j'ai contribué à des activités de développement logiciel et de coordination technique pour le segment sol de la mission spatiale \textbf{LISA}, dans un environnement CNES/ESA associant chaînes de données mission, niveaux L0--L3, cycle en V, intégration de modèles, validation et revue collaborative. J'y ai développé \textbf{CATpyLISA}, une plateforme Python destinée à structurer des workflows, comparer des modèles, qualifier des écarts, valider des résultats et produire une documentation exploitable par différents contributeurs.

Cette expérience est directement transposable à un rôle de chef de projet technique : qualifier une demande, clarifier les besoins, identifier les dépendances, rendre visibles les risques, organiser les actions, produire des synthèses, préparer les arbitrages et assurer le transfert de connaissances vers les équipes qui prennent le relais. Elle m'a également donné une pratique concrète du cycle en V, de l'IVVQ, de la traçabilité, de la documentation et de la coordination entre profils scientifiques, logiciels et institutionnels.

Plus récemment, chez \textbf{VEV Platform Services France}, j'ai travaillé sur une plateforme logicielle industrielle connectée à des infrastructures de recharge électrique, avec des flux de données opérationnels, des contraintes terrain, du support technique, de l'analyse d'incidents et de la documentation. Cette expérience m'a permis de renforcer ma compréhension des enjeux run : continuité de service, analyse de comportements applicatifs, diagnostic, formalisation des constats et contribution à la fiabilisation de services numériques.

Pour la CNAF, je peux apporter un profil capable de jouer un rôle de \textbf{SPOC technique} entre différentes parties prenantes : écouter les besoins, consolider les contraintes, suivre les actions, identifier les risques, alimenter les comités, produire un reporting fiable et maintenir un bon niveau d'information entre les équipes. Je suis habitué à travailler avec rigueur, méthode et autonomie sur plusieurs sujets en parallèle, tout en conservant une communication claire et structurée.

Je dispose également d'une culture solide des environnements SI et logiciels : Python, TypeScript, Angular, Node.js, Linux, Git, Docker, CI/CD, documentation technique, monitoring/logs, OpenSearch, Sentry, Confluence, Jira et Zenhub. Sans me présenter comme expert infrastructure DSI, je possède un socle technique suffisant pour dialoguer efficacement avec des architectes, des équipes build, run, production, support et projet.

Le rôle proposé par la CNAF représente pour moi une continuité logique vers un poste de pilotage technique stratégique, au service d'une DSI à fort impact public. Je suis particulièrement motivé par la dimension transverse du poste, la nécessité de fédérer des acteurs différents, la qualification des demandes techniques, l'identification des risques et la contribution à une plateforme technique au service des équipes de développement.

Je serais heureux de pouvoir échanger avec vous afin de vous présenter plus en détail mon parcours et la manière dont mon expérience en coordination technique, validation, documentation, analyse de risques et interfaces SI peut contribuer efficacement aux missions de la Direction Technique de la CNAF.

Je vous prie d'agréer, Madame, Monsieur, l'expression de mes salutations distinguées.

\textbf{Guillaume Boileau}

}

\end{document}
